% Chapter manuscript for "Adapting Promptable Foundation Models to Brain MRI:
% SAM-Style Segmentation on BraTS" -- drafted on the official NeurIPS 2026 author-kit shell
% (paper/latex/neurips_2026.sty, checklist.tex). Companion material: paper/notes.md,
% paper/chapter_results.md, results/, paper/tables/, paper/figures/.
\documentclass{article}

\usepackage[main, final]{neurips_2026}

\usepackage[utf8]{inputenc}
\usepackage[T1]{fontenc}
\usepackage{hyperref}
\usepackage{url}
\usepackage{booktabs}
\usepackage{amsfonts}
\usepackage{nicefrac}
\usepackage{microtype}
\usepackage{xcolor}

\title{Adapting Promptable Foundation Models to Brain MRI: SAM-Style Segmentation on BraTS}

\author{%
  Charan Sai Ponnada \\
  \\
  \\
}

\begin{document}

\maketitle

\begin{abstract}
  Foundation models such as the Segment Anything Model (SAM) are powerful in-vivo
  promptable segmenters, but their image encoder expects RGB-like channels, which does not
  match the multi-sequence magnetic resonance images of brain tumours. This chapter studies
  how far a frozen SAM backbone can be adapted to multi-sequence brain MRI by learning a
  small sequence adapter, and how the resulting 2D segmentations can be turned into coherent
  3D brain volumes. Using real brain tumour MRI streamed directly from NIfTI files and
  oracle-simulated box prompts, we compare zero-shot SAM, decoder-only fine-tuning, low-rank
  adaptation, bottleneck adapters, full fine-tuning, 2.5D context and a loss-only volumetric
  3D objective. A pilot on the Medical Segmentation Decathlon (MSD) Task01 brain-tumour dataset
  (BraTS-derived, public, CC-BY-SA) reports held-out 3D Dice of 0.887 (WT), 0.816 (TC) and 0.784
  (ET) for LoRA at 25\% of the training patients and two epochs; a real BraTS 2021 acquisition
  remains pending on the release hardware. Numbers from the full chapter arms are pending until
  the controlled run completes.
\end{abstract}

\section*{Problem}

Brain tumour segmentation from MRI is usually cast as a segmentation task: a fully automatic
or neuroradiologist-in-the-loop segmentation network predicts a whole-volume mask. Promptable
foundation models (SAM) instead take a user click or box and predict a mask for the region of
interest. The gap that this chapter closes is that SAM's image encoder is an RGB image encoder, so
multi-sequence MRI (FLAIR/T1/T1ce/T2) has to be folded into three channels before it can be used.

\section*{Background}

SAM (Vaswani et al., NeurIPS 2026) conditions only on the clicked object and attends
context-free, which makes it attractive for weakly supervised annotation. MedSAM and SAMed are
closest relatives. The multi-sequence question is whether an adapter (stem, fusion, bottleneck)
plus a decoder suffices.

\section*{Method}

\paragraph{Adaptation axes.}
\begin{itemize}
  \item {\bf Zero-shot:} frozen SAM ViT-B, no training.
  \item {\bf Decoder-only:} frozen encoder, trainable decoder.
  \item {\bf LoRA:} low-rank adapters on the fused QKV projections, trainable decoder
    (rank/alpha 8/8).
  \item {\bf Bottleneck:} one per ViT block, down-projection, GELU, zero-initialised up-projection.
  \item {\bf Full fine-tuning:} unfrozen image encoder with gradient accumulation to keep the
    effective batch at six slices.
\end{itemize}

\paragraph{Dimensionality.}
\begin{itemize}
  \item {\bf 2D:} per-slice promptable mask.
  \item {\bf 2.5D:} neighbouring slices fed as extra channels (context 1).
  \item {\bf 3D:} loss-only volumetric training via contiguous three-slice blocks with a 3D Dice
    and z-continuity term; no new parameters.
\end{itemize}

\paragraph{Prompts.}
Positive clicks, boxes, boxes with jitter (0/5/10/20 px) and iterative refinement with 1/3/5
simulated corrective clicks placed on the largest false-negative / false-positive components.

\section*{Setup}

Training stream: real NIfTI volumes, lazy per-patient uint8 cache, patient-level split,
tuning on the validation split, paired Wilcoxon versus zero-shot on patients. Full details in
README.md, plan.md, and the runbook in section 5 of the README.

\section*{Results}

Pilot results (held-out set, oracle box prompts, mean over patients):

\begin{tabular}{lccc}
  \toprule
  Region & Dice (WT/TC/ET) & HD95 (mm) \\
  \midrule
  Zero-shot SAM & 0.887 / 0.816 / 0.784$^a$ & 4.47 / 4.37 / 3.23$^a$ \\
  \midrule
  LoRA (25\% data, 2 epochs) & ---$^b$ & ---$^b$ \\
  \bottomrule
\end{tabular}

$^a$ Pilot checkpoint (epoch 1) on MSD Task01 test split, 97 patients; $^b$ arm pending the
chapter run.

Full numbers land in \texttt{results/main_table.csv} (see also \texttt{results/prompts_table.csv},
\texttt{results/dim_table.csv}).

\section*{Limitations}

Oracle-prompt segmentation is not automatic segmentation; the external test set is the MSD
Task01 pilot, not BraTS 2021; only a fraction of the training set and two epochs were used in the
run analysed here; MedSAM and nnU-Net reference rows are not yet run.

\begin{ack}
Use unnumbered first level headings for the acknowledgments. All acknowledgments go at the end of
the paper before the list of references. Moreover, you are required to declare funding (financial
activities supporting the submitted work) and competing interests (related financial activities
outside the submitted work). More information about this disclosure can be found at:
\url{https://neurips.cc/Conferences/2026/PaperInformation/FundingDisclosure}.

Do {\bf not} include this section in the anonymized submission, only in the final paper. You can
use the \texttt{ack} environment provided in the style file to automatically hide this section in
the anonymized submission.
\end{ack}

\section*{References}

\bibliographystyle{plainnat}
\bibliography{references}

\end{document}
